\chapter{Method}
\label{ch:method}

Training reuses the ShapeLLM PPO loop
(\texttt{environment.outer\_rollout}, \texttt{PPOAgent},
\texttt{CustomPPOTrainer}) on the logistic CPR of Chapter~\ref{ch:env}
\cite{garciasegura2025}.
This chapter defines the policy, the reward handed to the optimiser, the
advantage estimator and its two normalisation operators, the PPO objective
with every hyperparameter, and the two training schedules, naive and
shaper, whose difference is the object of Chapter~\ref{ch:results}.
The scientific object is that game's reward, not a new optimiser.

\section{Model, tokens, generation}

Base model: \texttt{google/gemma-2-2b-it}, \texttt{torch.bfloat16}, eager
attention.
Actions are single digit tokens (Table~\ref{tab:tokens}).

\begin{table}[ht]
\centering
\caption{Legal action tokens for Gemma-2-2b-it.
  \texttt{AllowedTokensLogitsProcessor} masks every other vocabulary
  entry.}
\label{tab:tokens}
\begin{tabular}{rr}
\toprule
Action & Token id \\
\midrule
0 & 235276 \\
1 & 235274 \\
2 & 235284 \\
3 & 235304 \\
\bottomrule
\end{tabular}
\end{table}

\paragraph{Policy.}
Each learning agent carries a rank-2 LoRA adapter \(\theta\) on the query
and value projections \cite{hu2022} and a scalar value head
\(V_{\theta}\) on the final hidden state (TRL's
\texttt{AutoModelForCausalLMWithValueHead}); the base weights are frozen.
At round \(t\) the agent receives a prompt \(x_t\) rendered from its
observable history (Chapter~\ref{ch:env}): the rules, the current stock
and, from round~2, the previous round's requests and receipts.
A shaper with \texttt{transmit\_info} also sees the joint-request counts
and episode summaries of the current trial.
The action is one token drawn from the softmax over the four legal action
tokens, every other vocabulary entry being floored by the logits
processor:
\begin{equation}
  \pi_{\theta}(a\mid x_t)
  =
  \frac{\exp z_{\theta}(x_t)_a}{\sum_{a'\in\mathcal{A}}\exp z_{\theta}(x_t)_{a'}},
  \qquad a\in\mathcal{A},
  \label{eq:policy}
\end{equation}
where \(z_{\theta}(x_t)\) are the logits at the answer position.
Generation is one new token (\texttt{max\_new\_tokens=1},
\texttt{do\_sample=True}, \texttt{top\_p=1.0}, \texttt{top\_k=0},
temperature~1), so the emitted token is exactly a draw from
\eqref{eq:policy}.
Illegal draws (observed under MPS multinomial, not under the mask) are
rejection-sampled up to four times.

\paragraph{Untrained prior.}
The pre-flight script (\texttt{scripts/cpr\_preflight.py}) measures the
untrained distribution at about \(0.90\) on token~2.
At the opening prompt of the executed runs, the first episode of every run
(three games before any update) opened 2 on 205 of 240 games, 1 on 25 and
3 on 10, so the untrained opening mass is about \(0.85\) on token~2,
\(0.10\) on token~1, \(0.04\) on token~3 and effectively zero on token~0
(these draws are not independent; see the seeds paragraph below).
That is an executed observation, not a literature citation.
The method does \textbf{not} flatten the prior, does \textbf{not} turn on
\texttt{use\_score\_scaling} / \texttt{use\_score\_norm}, and does
\textbf{not} add a take-1 bonus.
An entropy coefficient of \(0.15\) in early runs bought take-3, not
restraint; all live configurations use \(0.05\), decaying to \(0.01\).

\section{Learning algorithm}
\label{sec:algorithm}

\paragraph{Per-step reward.}
The reward handed to the optimiser at round \(t\) is the receipt
\(c_t\in\{0,1,2,3\}\) of \eqref{eq:harvest}, penalised by the divergence of
the sampled token from the frozen reference model, which is the same
network with the adapter disabled:
\begin{equation}
  \tilde{r}_t
  =
  c_t-\beta_t\bigl[\log\pi_{\theta}(a_t\mid x_t)-\log\pi_{\mathrm{ref}}(a_t\mid x_t)\bigr].
  \label{eq:klreward}
\end{equation}
The coefficient \(\beta_t\) follows the adaptive controller of Ziegler et
al.\ \cite{ziegler2019} as implemented in TRL~0.11.4: initial value
\(0.2\), target divergence \(6\), horizon \(10^{4}\) steps.
With batches of at most \(108\) positions the controller barely moves; the
logged value falls from \(0.200\) to \(0.181\) over the 75 updates of a
15-epoch naive run.
The term is zero under the untrained prior and grows as the policy leaves
it: at \(\pi_{\theta}(1\mid x_1)=0.9\) against \(\pi_{\mathrm{ref}}(1\mid
x_1)=0.1\) it is \(0.2\times\log 9\approx 0.44\) per step, comparable with
the one-unit receipt difference between neighbouring actions.
It is a restoring force towards the prior and is part of every result in
Chapter~\ref{ch:results}.
Steps after collapse (\(R_t=0\)) carry a NaN reward and are dropped from
the batch before any statistic is computed; the collapse step itself is a
real decision and stays.
\texttt{game.records} stores the true zero for those steps.

\paragraph{Advantages.}
Advantages are generalised advantage estimates \cite{schulman2016}
computed by a backward pass over each parallel game,
\begin{equation}
  \delta_t=\tilde{r}_t+\gamma V_{\theta}(x_{t+1})-V_{\theta}(x_t),
  \qquad
  \hat{A}_t=\sum_{l\ge 0}(\gamma\lambda)^{l}\,\delta_{t+l},
  \qquad \gamma=1,\ \lambda=0.97,
  \label{eq:gae}
\end{equation}
where the pass runs over one episode for a naive learner and over the
concatenated episodes of a trial for a shaper, so that with \(\gamma=1\) a
shaper's opening-step advantage includes receipts from later episodes of
the same trial.
The value target is the GAE return \(\hat{G}_t=\hat{A}_t+V_{\theta}(x_t)\).
Before entering the loss, the advantages of a batch \(\mathcal{B}\) are
transformed by one of two operators, where \(\mu_{\mathcal{B}}\) and
\(\sigma_{\mathcal{B}}\) are the mean and standard deviation over the
unmasked positions of the batch:
\begin{equation}
  \mathcal{N}_{\mathrm{whiten}}(\hat{A}_t)
  =
  \frac{\hat{A}_t-\mu_{\mathcal{B}}}{\sqrt{\sigma_{\mathcal{B}}^2+\varepsilon}},
  \qquad
  \mathcal{N}_{\mathrm{centre}}(\hat{A}_t)
  =
  \hat{A}_t-\mu_{\mathcal{B}}.
  \label{eq:norm}
\end{equation}
Whitening is the TRL default (\texttt{masked\_whiten}, the OpenAI
\texttt{lm-human-preferences} convention documented by Huang et al.\
\cite{huang2024}); centring is the \texttt{advantage\_norm=center} flag of
this repository.
Masked positions are zeroed after either operator.
Engstrom et al.\ \cite{engstrom2020} show that PPO's measured behaviour is
driven by such code-level choices, not only by the clip.

\paragraph{Objective.}
With \(\tilde{A}_t=\mathcal{N}(\hat{A}_t)\), the probability ratio
\(\varrho_t=\pi_{\theta}(a_t\mid x_t)/\pi_{\theta_{\mathrm{old}}}(a_t\mid x_t)\),
the old value estimate \(V_{\mathrm{old}}\) and the clipped value
\(V^{\mathrm{c}}=\operatorname{clip}(V_{\theta},V_{\mathrm{old}}-\epsilon_v,V_{\mathrm{old}}+\epsilon_v)\),
each agent maximises the clipped surrogate of Schulman et al.\
\cite{schulman2017}
\begin{equation}
  \begin{aligned}
  \mathcal{L}(\theta)=\mathbb{E}_{\mathcal{B}}\Bigl[
    &\min\bigl(\varrho_t\tilde{A}_t,\;
      \operatorname{clip}(\varrho_t,1-\epsilon,1+\epsilon)\,\tilde{A}_t\bigr)\\
    &-\tfrac{c_v}{2}\max\bigl((V_{\theta}(x_t)-\hat{G}_t)^2,\,(V^{\mathrm{c}}(x_t)-\hat{G}_t)^2\bigr)
    +c_e\,\mathcal{H}\bigl(\pi_{\theta}(\cdot\mid x_t)\bigr)
  \Bigr],
  \end{aligned}
  \label{eq:ppo}
\end{equation}
where the entropy is taken over the four legal tokens only.
The policy clip is \(\epsilon=0.2\) for naive learners (the TRL default;
no live configuration overrides it) and \(0.1\) for the shaper and for
agent~2 of the slow-learning-rate control; the value clip is
\(\epsilon_v=0.2\); \(c_v=0.01\); and \(c_e\) starts at \(0.05\) and falls
by \(0.04/150\) per update with a floor of \(0.01\), so it reaches
\(0.03\) after the 75 updates of a naive run and \(0.046\) after the 15
updates of a shaper.
The optimiser is Adam \cite{kingma2015} with the learning rate of
Table~\ref{tab:hyperparameters} and default moments
(\(\beta_1=0.9\), \(\beta_2=0.999\), \(\varepsilon=10^{-8}\)); one pass
over the batch in minibatches of five; no gradient clipping.
The batch is the set of unmasked positions of one episode (naive) or one
trial (shaper), so its size varies with survival.

\begin{table}[ht]
\centering
\small
\caption{Every hyperparameter of the executed runs.
  ``Shaper'' is agent~2 of the naive--shaper configurations; ``control''
  is agent~2 of the slow-learning-rate naive--naive configuration, which
  copies the shaper's learning rate and clip but keeps the naive schedule
  and prompt.
  Agent~1 of every two-learner configuration, and the learner of Tests
  A~and~B, are naive learners.}
\label{tab:hyperparameters}
\begin{tabularx}{\textwidth}{>{\raggedright\arraybackslash}p{4.3cm}>{\raggedright\arraybackslash}X}
\toprule
Quantity & Value \\
\midrule
Base model & \texttt{gemma-2-2b-it}, bf16, eager attention, base weights frozen \\
Adapter & LoRA \(r=2\), \(\alpha=32\), dropout \(0.05\), on \texttt{q\_proj} and \texttt{v\_proj} \\
Value head & scalar, on the final hidden state \\
Generation & one token, sampling, \texttt{top\_p} 1, \texttt{top\_k} 0, temperature 1 \\
Learning rate (Adam) & naive \(1.41\times 10^{-6}\); shaper \(3\times 10^{-7}\); control \(3\times 10^{-7}\) \\
Adam moments & \(\beta_1=0.9\), \(\beta_2=0.999\), \(\varepsilon=10^{-8}\); no gradient clipping \\
Policy clip \(\epsilon\) & naive 0.2; shaper 0.1; control 0.1 \\
Value clip \(\epsilon_v\), coefficient \(c_v\) & 0.2, 0.01 \\
Entropy coefficient \(c_e\) & \(0.05\to 0.01\), \(-0.04/150\) per update, over the four legal tokens \\
KL coefficient \(\beta_t\) & adaptive, initial 0.2, target 6, horizon \(10^{4}\); logged \(0.200\to 0.181\) over 75 updates \\
\(\gamma\), \(\lambda\) & 1, 0.97 \\
GAE span & naive and control: one episode; shaper: one trial of 5 episodes \\
Update schedule & naive and control: after each episode; shaper: once per trial \\
Batch & unmasked positions of the span (\(\le 108\) per episode, \(\le 540\) per trial) \\
Minibatch, PPO epochs & 5, 1 \\
Trial prompt & shaper only: info-on (counts and episode summaries) or info-off \\
Advantage operator & centre in every two-learner and Stage~B run; whiten in the whitened Test~A/B runs \\
Trial geometry & \(T=36\) rounds, \(E=5\) episodes per trial, 3 parallel games, 15 games per epoch \\
Epochs & 15; 30 for seed-fixed whitened Test~A seed~0; 50 for the long seed-0 runs \\
Updates per 15 epochs & naive and control 75; shaper 15 \\
Seeds & 0, 1, 2 (see the seeds paragraph) \\
Hardware & NVIDIA L40S (RunPod) for Stage~B, the seed-fixed packet, whitened seeds 1--2 and the 50-epoch naive run; Apple silicon (MPS) for the pre-fix Test~A/B seed-0 tapes \\
Software & \texttt{trl} 0.11.4, \texttt{transformers} 4.47.0, \texttt{peft} \\
\bottomrule
\end{tabularx}
\end{table}

\paragraph{Seeds.}
Each experiment \(k\) calls \texttt{set\_seed}\((k)\) before constructing
its agents.
TRL's \texttt{PPOTrainer} then calls \texttt{set\_seed}\((0)\) in its own
constructor (its \texttt{PPOConfig} default), after the adapter and value
head have been created.
On every run dated 6~September or earlier the action sampler was therefore
drawn from the seed-0 stream, and the nominal seed controlled only the
value-head initialisation.
Those runs of nominally different seeds on the same hardware begin
identically and diverge only through non-deterministic bf16 kernels and,
in Stage~B, through the noise tables, which are seeded separately and
correctly.
\texttt{scripts/audit\_seeds\_and\_scale.py} measures the consequence: the
pre-fix whitened Test~A seeds 1 and 2 opened identically on all 75 games
of their first five epochs, and the Stage~B naive--naive, naive--shaper
and slow-LR seeds 1 and 2 opened identically on 92--95\% of games over
the whole run.
The entry points re-seed after construction.
The 8--9~September packet used that fix: seed-fixed Test~A centre pairs
agree on \(0.56\)--\(0.61\) of openings in epochs 1--5 (pre-fix whitened
\(s1\) vs \(s2\) was \(1.00\)); NS~\(\xi\) seed~0 is identical to the
pre-fix seed-0 tape, as expected.
Chapter~\ref{ch:lim} states what this does to the inference.
Write ``three seeds'' for the seed-fixed packet; write ``three nominal
seeds'' for every earlier GPU table.

\section{What the normalisation operator changes}
\label{sec:operator}

The opening-step advantage \(A_0\) is logged for the first index of each
parallel game in the flattened batch, raw and after the operator.
At \(\varrho_t=1\) the policy-gradient term of \eqref{eq:ppo} has expected
gradient on the logit of opening action \(k\)
\begin{equation}
  \frac{\partial\mathcal{L}}{\partial z_k}\Big|_{x_1}
  =
  \pi(k\mid x_1)\bigl(\tilde{A}(k)-\bar{A}\bigr),
  \qquad
  \bar{A}=\sum_{a\in\mathcal{A}}\pi(a\mid x_1)\,\tilde{A}(a),
  \label{eq:logit-update}
\end{equation}
where \(\tilde{A}(a)\) is the normalised opening advantage of action
\(a\), because \(\partial\log\pi(a\mid x)/\partial z_k=\mathbb{1}[a=k]-\pi(k\mid x)\)
for a softmax.
The sign of the update on logit \(k\) is the sign of
\(\tilde{A}(k)-\bar{A}\), so an action whose advantage is positive but
below the policy-weighted mean is pushed down; and the magnitude is
proportional to \(\pi(k\mid x_1)\), so a rare action moves slowly however
large its advantage.
Both operators in \eqref{eq:norm} are affine with positive scale, so they
agree on every sign; they differ in the scale of the gradient by
\(\sigma_{\mathcal{B}}\).

\paragraph{The measured scale.}
Per update, from the \texttt{live\_adv} logs of the six Test~A runs, the
batch standard deviation of the raw advantages has median \(13\)--\(15\)
(interquartile range \(12\)--\(15\) on the centred runs).
On the whitened runs it falls to \(2\)--\(5\) in the updates where every
game collapsed by round~4, because such a batch contains almost no
variation to divide by: whitening then scales an uninformative batch up to
unit variance.
An earlier estimate of \(53\), inferred from run-averaged opening
advantages, was not a per-batch quantity and is withdrawn.

\paragraph{Why the scale is not a step size.}
Adam's update \(\hat{m}/(\sqrt{\hat{v}}+\varepsilon)\) is invariant to a
uniform rescaling of the gradient \cite{kingma2015}, so a run in which
every batch were rescaled by the same \(\sigma_{\mathcal{B}}\) would take
the same steps under either operator.
Multiplying the learning rate by \(\sigma_{\mathcal{B}}\) does not convert
a whitened run into a centred one.
What the operator changes is the weight of the policy-gradient term
relative to two things that are not rescaled.
First, the value loss and the entropy bonus in \eqref{eq:ppo}: under
whitening the policy term is of order one per position while the value
regression on returns of \(10\)--\(100\) contributes a gradient of order
\(c_v\,|V_{\theta}-\hat{G}|\approx 0.01\times 40\) through the shared
adapter, so the two are comparable; under centring the policy term is
\(13\)--\(15\) times larger and dominates.
Second, the other batches: Adam's second moment is a running average over
updates, so the step taken on a batch scales with that batch's gradient
norm relative to the recent history.
Under whitening every batch has unit advantage variance, and a batch in
which every game collapsed drives a step of the same size as a batch
containing a surviving game; under centring the surviving batch carries a
gradient several times larger than a collapse-only batch
(\(\sigma_{\mathcal{B}}\approx 14\) against \(2\)) and takes a
proportionally larger step.
Centring therefore weights informative batches over uninformative ones;
whitening equalises them and raises the relative weight of the critic and
the entropy bonus.
Which of these produced the difference between the whitened and centred
seed-0 runs is not identified by the executed runs; the ablation that
would identify it is whitening with \(c_v\) and \(c_e\) divided by
\(\sigma_{\mathcal{B}}\), or plain SGD in place of Adam, neither of which
was run.

\paragraph{Worked example.}
With the first-episode opening distribution \(\pi(2)\approx 0.85\),
\(\pi(1)\approx 0.10\) and the centred seed-0 opening advantages of
Chapter~\ref{ch:results} (\(\tilde{A}(1)=+16.7\), \(\tilde{A}(2)=-5.7\),
\(\bar{A}\approx -3.2\)), \eqref{eq:logit-update} gives a gradient of
about \(+2.0\) on logit~1 and \(-2.1\) on logit~2.
Dividing by \(\sigma_{\mathcal{B}}\approx 13\) gives the whitened
counterparts \(+0.15\) and \(-0.16\) with the same signs.
The quantity that differed between the whitened seed-0 run and the
whitened seeds 1--2 is not this scale but the sign of
\(\tilde{A}(2)-\bar{A}\): in the seed-0 run the post-operator opening-2
advantage stayed between \(+0.5\) and \(+1.1\) in every epoch because
nearly every game in every batch collapsed, so there was no surviving game
in the batch for the opening-2 games to be compared against, whereas in
seeds 1--2 it turned negative from epoch~4, once surviving games entered
the batches (Table~\ref{tab:a0-epochs}).

\section{Frozen-bot tests (executed path)}

Before any shaper-versus-naive grid, one PPO learner is trained against a
non-learning partner.
This is a \emph{precondition} measurement, not opponent shaping: if the
prior never leaves opening 2 against a partner that is already frozen
on 2, a later two-learner table in which one agent leaves 2 cannot be
read as ``the shaper taught restraint.''
The two-learner and Stage~B grids therefore use mean-centring
(Table~\ref{tab:hyperparameters}), the operator under which leave-2 is
reliable within 15 epochs on this lock.
The TRL default (whitening) is retained as the Test~A/B contrast, not
hidden.

\begin{itemize}
\item \textbf{Test A.} \texttt{ConstantActionAgent} always requests 2
  (frozen hawk).
\item \textbf{Test B.} Same bot always requests 1 (frozen dove).
\end{itemize}

The bot has no Gemma weights.
\texttt{is\_shaper=False} so \texttt{outer\_rollout} calls a no-op
\texttt{update\_parameters} once per episode.
Only the learner's LoRA and value head train.
Records still include both request columns.

Trial geometry in the Test~A/B configs: \(T=36\), \(E=5\)
episodes per trial, \(3\) parallel games, 15 epochs.
That is 15 games per epoch and 225 games per run
(Figure~\ref{fig:trial}).

\begin{figure}[ht]
\centering
\begin{tikzpicture}[
  font=\small,
  epibox/.style={
    draw, rounded corners=2pt, minimum width=1.55cm,
    minimum height=0.62cm, inner sep=1pt,
  },
  arr/.style={-{Latex}, thick},
  player/.style={circle, inner sep=2.0pt},
]
\node[font=\scriptsize] at (4.2, 3.05)
  {one epoch $=$ one trial of \(3\times 5=15\) games};
\foreach \e/\x in {1/0, 2/2.1, 3/4.2, 4/6.3, 5/8.4} {
  \node[font=\scriptsize] at (\x, 2.55) {$e=\e$};
  \node[epibox, fill=blue!32] (n1e\e) at (\x, 1.70) {};
  \node[epibox, fill=orange!38] (n2e\e) at (\x, 0.85) {};
  \node[epibox, fill=black!12] (n3e\e) at (\x, 0.00) {};
}
\foreach \e in {1,2,3,4} {
  \pgfmathtruncatemacro{\nxt}{\e+1}
  \draw[arr] (n1e\e) -- (n1e\nxt);
  \draw[arr] (n2e\e) -- (n2e\nxt);
  \draw[arr] (n3e\e) -- (n3e\nxt);
}
\node[left=0.18cm of n1e1, font=\scriptsize] {$n=1$};
\node[left=0.18cm of n2e1, font=\scriptsize] {$n=2$};
\node[left=0.18cm of n3e1, font=\scriptsize] {$n=3$};
\draw[decorate, decoration={brace, amplitude=4pt, mirror}]
  ([yshift=-5pt]n3e3.south west) -- ([yshift=-5pt]n3e3.south east)
  node[midway, below=4pt, font=\scriptsize] {$T=36$ rounds};
% one episode, round by round
\node[player, fill=gray!55] (p1a) at (0.55, -1.55) {};
\node[player, fill=black, right=3pt of p1a] (p2a) {};
\node[font=\scriptsize, above=1pt of p1a, xshift=5pt] {$t=1$};
\node[player, fill=gray!55] (p1b) at (2.65, -1.55) {};
\node[player, fill=black, right=3pt of p1b] (p2b) {};
\node[font=\scriptsize, above=1pt of p1b, xshift=5pt] {$t=2$};
\node[font=\scriptsize] (dots) at (4.55, -1.55) {$\cdots$};
\node[player, fill=gray!55] (p1c) at (6.15, -1.55) {};
\node[player, fill=black, right=3pt of p1c] (p2c) {};
\node[font=\scriptsize, above=1pt of p1c, xshift=5pt] {$t=T$};
\draw[arr] (p2a) -- (p1b);
\draw[arr] (p2b) -- (dots);
\draw[arr] (dots) -- (p1c);
\draw[dashed, rounded corners=3pt]
  (-0.35, -2.05) rectangle (8.95, -1.05);
\draw[dashed, -{Latex}] (n3e1.south) -- ++(0,-0.35) -| (0.55, -1.05);
\node[font=\scriptsize, anchor=north] at (4.2, -2.20)
  {naive: PPO after each episode\qquad shaper: PPO once per trial};
\end{tikzpicture}
\caption{One training epoch of the CPR loop, drawn after Garcia Segura
  et al.\ (their Figure~1).
  Rows are three independent parallel copies of the same game.
  Columns are five sequential episodes of \(T=36\) rounds.
  Same-coloured boxes are successive episodes of one parallel copy.
  The dashed box is one episode: both agents act simultaneously each
  round.
  A naive learner updates after every episode; a shaper updates once
  from the whole trial.
  Fifteen such epochs are 225 games.
  Frozen-bot Test~A/B uses the same geometry; the partner does not
  update.}
\label{fig:trial}
\end{figure}

Versus always-2, a constant take-1 scores 36 and lives; constant take-2
scores 7 and dies.
Versus always-1, constant take-2 scores 72 and lives; constant take-1
scores 36; take-3 dies.
Test~B does not treat ``copy 1'' as the unique high-return reply.

\section{Naive and shaping schedules}
\label{sec:schedules}

\texttt{finetuning\_cpr.py} trains two \texttt{PPOAgent}s through the same
\texttt{outer\_rollout}.
Training proceeds in trials of \(E=5\) episodes
(Figure~\ref{fig:trial}).
A naive learner updates its parameters after every episode \(e\) from
that episode's batch,
\begin{equation}
  \theta_{e+1}
  =
  \theta_e+\alpha_{\mathrm{n}}\,\nabla_{\theta}
  \mathcal{L}\bigl(\theta_e;\ \tau_e(\pi_{\theta_e},\pi_{\phi})\bigr),
  \label{eq:naive-update}
\end{equation}
where \(\tau_e\) is the trajectory of episode \(e\) and \(\pi_{\phi}\) is
the partner (the gradient step is Adam's, written here as a plain step for
readability).
A shaper holds \(\phi\) fixed for the whole trial and updates once, from
the concatenated trajectories, towards the model-free opponent-shaping
objective of Lu et al.\ \cite{lu2022} as instantiated for language-model
agents by Garcia Segura et al.\ \cite{garciasegura2025}:
\begin{equation}
  J(\phi)
  =
  \mathbb{E}\Bigl[
    \sum_{e=1}^{E}
    G_{\mathrm{shaper}}^{(e)}\bigl(\pi_{\phi},\pi_{\theta_e}\bigr)
  \Bigr],
  \qquad
  \theta_{e+1}=\theta_e+\alpha_{\mathrm{n}}\nabla_{\theta}
  \mathcal{L}\bigl(\theta_e;\tau_e(\pi_{\theta_e},\pi_{\phi})\bigr).
  \label{eq:shaper-objective}
\end{equation}
The naive objective \eqref{eq:naive-update} treats the partner as fixed
within the episode; the shaper objective \eqref{eq:shaper-objective}
credits \(\phi\) for the partner's parameter change
\(\theta_e\to\theta_{e+1}\) that \(\phi\) induced, through the trial-level
GAE of \eqref{eq:gae}.
That dependence is the only thing that distinguishes shaping from
ordinary learning against a slowly changing opponent, and it can only be
exploited if the shaper's own policy changes on a timescale at which the
naive learner responds.
The shaper's learning rate \(\alpha_{\mathrm{s}}=3\times 10^{-7}\),
tighter clip and once-per-trial schedule therefore also make it a slower,
more stationary opponent than a matched-learning-rate naive learner.
The slow-learning-rate naive--naive control of
Section~\ref{sec:executed-two} gives agent~2 the same learning rate and
clip but the naive schedule and no trial prompt, and so isolates that side
effect from the shaping term.
When \texttt{transmit\_info=true}, the shaper prompt also keeps
joint-request counts and episode summaries for the trial.
Live protocol:
\path{cpr_naive_naive_center.json},
\path{cpr_naive_shaper_center.json},
\path{cpr_naive_shaper_center_info_off.json},
\path{cpr_naive_naive_slow2_center.json}, and their \texttt{\_xi}
counterparts.

\texttt{finetuning\_fixed\_opponent.py} is not used for CPR.

\section{Stage B noise (machinery)}

Stage~B configs (\texttt{cpr\_*\_xi.json}) attach a common-random-number
table of \(\texttt{xi\_tenths}\in\{7,10,13\}\) per seed
(\texttt{cpr\_env.make\_noise\_table}).
The same table is shared across arms so arm contrasts are not confounded
by different draws.
\(\texttt{xi\_tenths}=10\) is bit-identical to Stage~A.
Test~A \(\xi\) uses \texttt{finetuning\_cpr\_fixed.py}; two-learner
\(\xi\) uses \texttt{finetuning\_cpr.py}.
Tables are saved as \texttt{noise\_table.npy} next to the records.
This is a measurement device, not a new optimiser.
